% !TEX program = pdflatex
\documentclass[11pt]{article}

% ==========================
% Encoding / font (pdflatex)
% ==========================
\usepackage[utf8]{inputenc}
\usepackage[T1]{fontenc}
\usepackage[english]{babel}
\usepackage{lmodern}
\usepackage{microtype}

% ==========================
% Page layout
% ==========================
\usepackage[margin=1in]{geometry}

% ==========================
% Math
% ==========================
\usepackage{amsmath,amssymb,amsthm,mathtools}

% ==========================
% Lists
% ==========================
\usepackage{enumitem}

% ==========================
% Hyperref (load late)
% ==========================
\usepackage[hidelinks]{hyperref}

% ==========================
% Macros
% ==========================
\newcommand{\Pbb}{\mathbb{P}}
\newcommand{\E}{\mathbb{E}}
\newcommand{\R}{\mathbb{R}}
\newcommand{\F}{\mathcal{F}}
\newcommand{\1}{\mathbf{1}}
\newcommand{\dd}{\,\mathrm{d}}

% ==========================
% Theorem environments
% ==========================
\numberwithin{equation}{section}

\theoremstyle{plain}
\newtheorem{theorem}{Theorem}[section]
\newtheorem{lemma}[theorem]{Lemma}
\newtheorem{proposition}[theorem]{Proposition}
\newtheorem{corollary}[theorem]{Corollary}

\theoremstyle{definition}
\newtheorem{definition}[theorem]{Definition}

\theoremstyle{remark}
\newtheorem{remark}[theorem]{Remark}

% ==========================
% "Board" environment (robust)
%   - no extra spaces (ignorespaces/unskip)
%   - boxed minipage
% ==========================
\newenvironment{board}{%
  \par\medskip\noindent
  \fbox{%
    \begin{minipage}{0.98\linewidth}
      \textbf{On the board (plan + comments).}\par\medskip
      \ignorespaces
}{%
      \unskip
    \end{minipage}%
  }%
  \par\medskip
}

% ==========================
% Title
% ==========================
\title{Board Notes --- Constructing the It\^o Integral}
\author{}
\date{}

\begin{document}
\maketitle
\vspace{-1em}

\noindent\textbf{Goal of the lecture.}
We construct $\int_0^T \phi_t\,\dd W_t$ \emph{without} using It\^o's formula,
starting from left-point sums and passing to the limit in $L^2$.
We will establish:
\begin{itemize}[leftmargin=2em]
  \item the \textbf{correct} convention is the left-point (non-anticipative) one,
  \item the \textbf{key property} is the It\^o isometry,
  \item the integral is a \textbf{martingale} and we can compute its moments/covariances.
\end{itemize}

% ==========================================================
\section{Why Riemann--Stieltjes fails for Brownian motion}
% ==========================================================

\subsection{Infinite variation and nowhere differentiability}

\begin{remark}
If $x:[0,T]\to\R$ has \emph{finite total variation} on $[0,T]$, then the
Riemann--Stieltjes integral $\int_0^T \phi(t)\,\dd x(t)$ exists for every continuous $\phi$.
Brownian motion $W$ has (almost surely) infinite variation, which is a first reason for the failure.
\end{remark}


\begin{itemize}
  \item Recall total variation:
  \[
    \mathrm{Var}_{[0,T]}(x)=\sup_{\pi}\sum_{k}|x(t_{k+1})-x(t_k)|.
  \]
  \item Intuition: $\Delta W \sim \sqrt{\Delta t}$, so
  \[
    \sum |\Delta W|
    \sim \sum \sqrt{\Delta t}
    \approx \frac{T}{\sqrt{\Delta t}}\to\infty.
  \]
  \item Transition: rather than order-1 variation, what \emph{accumulates} for $W$ is order 2:
  \[
    \sum (\Delta W)^2.
  \]
\end{itemize}

\subsection{Quadratic variation: accumulation of the order-2 terms}

Fix a partition $\pi=\{0=t_0<\cdots<t_n=T\}$ and define
\[
Q_T^\pi(W)=\sum_{k=0}^{n-1}\big(W_{t_{k+1}}-W_{t_k}\big)^2.
\]

\begin{theorem}[Quadratic variation of Brownian motion: uniform case]
For the uniform partition $t_k=kT/n$,
\[
Q_T^{(n)}(W):=\sum_{k=0}^{n-1}(W_{t_{k+1}}-W_{t_k})^2 \xrightarrow[n\to\infty]{L^2} T.
\]
In particular, the convergence holds in probability.
\end{theorem}

\begin{proof}
Let $\Delta W_k:=W_{t_{k+1}}-W_{t_k}$. Then $\Delta W_k\sim\mathcal N(0,\Delta t)$,
independent, with $\Delta t=T/n$.

\smallskip
\noindent\textbf{(1) Expectation.}
\[
\E\big[Q_T^{(n)}(W)\big]=\sum_{k=0}^{n-1}\E[(\Delta W_k)^2]=\sum_{k=0}^{n-1}\Delta t = n\Delta t = T.
\]

\smallskip
\noindent\textbf{(2) Variance.}
Since the $\Delta W_k$ are independent,
\[
\mathrm{Var}\big(Q_T^{(n)}(W)\big)
=\sum_{k=0}^{n-1}\mathrm{Var}\big((\Delta W_k)^2\big).
\]
If $Z\sim\mathcal N(0,\sigma^2)$, then $\E[Z^4]=3\sigma^4$, hence
$\mathrm{Var}(Z^2)=\E[Z^4]-\E[Z^2]^2=3\sigma^4-\sigma^4=2\sigma^4$.
Here $\sigma^2=\Delta t$, so $\mathrm{Var}((\Delta W_k)^2)=2(\Delta t)^2$ and
\[
\mathrm{Var}\big(Q_T^{(n)}(W)\big)=n\cdot 2(\Delta t)^2
=2n\cdot (T/n)^2=\frac{2T^2}{n}\xrightarrow[n\to\infty]{}0.
\]
Therefore $\E\big[(Q_T^{(n)}(W)-T)^2\big]=\mathrm{Var}(Q_T^{(n)}(W))\to 0$, i.e.\ $L^2$-convergence.
\end{proof}

\textbf{To do on board}
\begin{itemize}
  \item Do this computation on the board (a great ``proof by computation'').
  \item Emphasize: this is \emph{the} crucial difference with finite-variation functions.
\end{itemize}
\textbf{To do on board}

\subsection{Why the Riemann--Stieltjes limit can depend on the sampling point}

Consider the random integrand $\phi_t=W_t$ and the sums
\[
L_n=\sum_{k=0}^{n-1}W_{t_k}\Delta W_k,
\qquad
R_n=\sum_{k=0}^{n-1}W_{t_{k+1}}\Delta W_k.
\]
Then
\[
R_n-L_n=\sum_{k=0}^{n-1}(W_{t_{k+1}}-W_{t_k})^2
=Q_T^{(n)}(W)\xrightarrow[]{\Pbb}T.
\]
Hence, \emph{if} $L_n$ converges, then $R_n$ converges to a different limit (shifted by $T$).
This shows cleanly that the usual Riemann--Stieltjes notion is not robust here.

\textbf{To do on board}
\begin{itemize}
  \item Say it very clearly: \textbf{changing the evaluation point changes the limit}
  when the integrand is as rough as $W$.
  \item Motivation: we need a new definition, and finance enforces the left-point choice
  (predictability / no peeking).
\end{itemize}
\textbf{To do on board}

% ==========================================================
\section{Finance interpretation: self-financing gains}
% ==========================================================

\subsection{Discrete $\to$ continuous: why the left-point convention}

On a grid $0=t_0<\cdots<t_n=T$, holding $\phi_{t_k}$ shares on $(t_k,t_{k+1}]$ yields the gain
\[
G_T^\pi(\phi)=\sum_{k=0}^{n-1}\phi_{t_k}\big(S_{t_{k+1}}-S_{t_k}\big).
\]
The key point is: \textbf{$\phi_{t_k}$ is chosen using the information at time $t_k$ (not after).}

\textbf{To do on board}
\begin{itemize}
  \item Direct link to ``no peeking'': $\phi_{t_k}$ is $\F_{t_k}$-measurable.
  \item This is the pedagogical reason for the left-point rule in the It\^o integral.
\end{itemize}
\textbf{To do on board}

% ==========================================================
\section{Setup: filtration, adaptedness, square integrability}
% ==========================================================

We work on $(\Omega,\F,(\F_t)_{t\ge0},\Pbb)$ supporting a Brownian motion $W$.
Typically $\F_t=\sigma(W_s:0\le s\le t)$ (completed).

\begin{definition}[Adapted / simple process]
A process $(\phi_t)$ is \emph{adapted} if $\phi_t$ is $\F_t$-measurable for every $t$.
A \emph{simple process} on $[0,T]$ is of the form
\[
\phi_t=\sum_{k=0}^{n-1}\phi_{t_k}\,\1_{(t_k,t_{k+1}]}(t),
\qquad \phi_{t_k}\in\F_{t_k}.
\]
\end{definition}

\begin{remark}[Why $L^2$?]
The natural condition to define $\int_0^T \phi_t\,\dd W_t$ is
\[
\E\!\left[\int_0^T \phi_t^2\,\dd t\right]<\infty,
\]
because the integral will have variance of that order (It\^o isometry).
\end{remark}

% ==========================================================
\section{Step 1: definition on simple processes (left sums)}
% ==========================================================

\subsection{Definition}

\begin{definition}[It\^o integral for simple processes]
For $\phi$ simple as above, define
\[
I_0(\phi):=\int_0^T \phi_t\,\dd W_t
:=\sum_{k=0}^{n-1}\phi_{t_k}\big(W_{t_{k+1}}-W_{t_k}\big).
\]
\end{definition}

\textbf{To do on board}
\begin{itemize}
  \item Stress: this is not a limit here; it is a \textbf{definition}.
  \item Stress: $\phi_{t_k}$ is evaluated \textbf{before} the increment $\Delta W_k$.
\end{itemize}
\textbf{To do on board}

\subsection{First quick properties}
\begin{proposition}[Linearity and measurability]
If $\phi,\psi$ are simple and $\alpha,\beta\in\R$, then
$I_0(\alpha\phi+\beta\psi)=\alpha I_0(\phi)+\beta I_0(\psi)$.
Moreover $I_0(\phi)$ is $\F_T$-measurable.
\end{proposition}

\begin{proof}
Linearity is immediate from linearity of sums.
For measurability: each $\phi_{t_k}$ is $\F_{t_k}\subset\F_T$-measurable and each increment
$W_{t_{k+1}}-W_{t_k}$ is $\F_{t_{k+1}}\subset\F_T$-measurable, hence their product is $\F_T$-measurable.
\end{proof}

% ==========================================================
\section{Step 2: the central computation --- mean zero and It\^o isometry}
% ==========================================================

\subsection{Mean}
\begin{proposition}
For $\phi$ simple, $\E[I_0(\phi)]=0$.
\end{proposition}

\begin{proof}
\[
\E[I_0(\phi)]=\sum_{k=0}^{n-1}\E\!\left[\phi_{t_k}(W_{t_{k+1}}-W_{t_k})\right]
=\sum_{k=0}^{n-1}\E\!\left[\E\!\left[\phi_{t_k}(W_{t_{k+1}}-W_{t_k})\mid \F_{t_k}\right]\right].
\]
Since $\phi_{t_k}$ is $\F_{t_k}$-measurable and the increment $(W_{t_{k+1}}-W_{t_k})$
is independent of $\F_{t_k}$ with mean $0$,
\[
\E\!\left[\phi_{t_k}(W_{t_{k+1}}-W_{t_k})\mid \F_{t_k}\right]
=\phi_{t_k}\,\E[W_{t_{k+1}}-W_{t_k}]=0.
\]
\end{proof}

\textbf{To do on board}
Do the conditioning carefully once: it is a model argument students should learn.
\textbf{To do on board}

\subsection{It\^o isometry (detailed proof for simple processes)}
\begin{theorem}[It\^o isometry for simple processes]
For $\phi$ simple,
\[
\E\!\left[I_0(\phi)^2\right]=\E\!\left[\int_0^T \phi_t^2\,\dd t\right]
=\sum_{k=0}^{n-1}\E[\phi_{t_k}^2]\,(t_{k+1}-t_k).
\]
\end{theorem}

\begin{proof}
Let $\Delta W_k=W_{t_{k+1}}-W_{t_k}$. Then $I_0(\phi)=\sum_k \phi_{t_k}\Delta W_k$ and
\[
I_0(\phi)^2=\sum_k \phi_{t_k}^2(\Delta W_k)^2
+2\sum_{i<j}\phi_{t_i}\phi_{t_j}\Delta W_i\Delta W_j.
\]
Taking expectations. For the diagonal term:
\[
\E\!\left[\phi_{t_k}^2(\Delta W_k)^2\right]
=\E\!\left[\E\!\left[\phi_{t_k}^2(\Delta W_k)^2\mid \F_{t_k}\right]\right]
=\E\!\left[\phi_{t_k}^2\,\E[(\Delta W_k)^2]\right]
=\E[\phi_{t_k}^2]\,(t_{k+1}-t_k),
\]
since $\Delta W_k$ is independent of $\F_{t_k}$ and $\E[(\Delta W_k)^2]=t_{k+1}-t_k$.

For cross terms $i<j$, condition on $\F_{t_j}$:
\[
\E[\phi_{t_i}\phi_{t_j}\Delta W_i\Delta W_j]
=\E\!\left[\E[\phi_{t_i}\phi_{t_j}\Delta W_i\Delta W_j\mid \F_{t_j}]\right].
\]
Now $\phi_{t_i},\phi_{t_j},\Delta W_i$ are $\F_{t_j}$-measurable (because $t_i<t_j$),
while $\Delta W_j$ is independent of $\F_{t_j}$ with mean $0$, hence
\[
\E[\phi_{t_i}\phi_{t_j}\Delta W_i\Delta W_j\mid \F_{t_j}]
=\phi_{t_i}\phi_{t_j}\Delta W_i\,\E[\Delta W_j]=0.
\]
So all cross terms vanish. Therefore
\[
\E[I_0(\phi)^2]=\sum_{k=0}^{n-1}\E[\phi_{t_k}^2]\,(t_{k+1}-t_k)
=\E\!\left[\int_0^T \phi_t^2\,\dd t\right].
\]
\end{proof}

\textbf{To do on board}
\begin{itemize}
  \item This is \textbf{the} computation of the lecture: take your time and show the structure.
  \item Message: \textbf{cross terms disappear} thanks to independence + mean zero.
\end{itemize}
\textbf{To do on board}

% ==========================================================
\section{Consequence: Cauchy property in $L^2(\Omega)$}
% ==========================================================

\begin{proposition}[Cauchy]
If $\phi^n$ and $\phi^m$ are simple processes, then
\[
\E\!\left[\big(I_0(\phi^n)-I_0(\phi^m)\big)^2\right]
=\E\!\left[\int_0^T (\phi_t^n-\phi_t^m)^2\,\dd t\right].
\]
Hence if $(\phi^n)$ is Cauchy in $L^2(\Omega\times[0,T])$, then $(I_0(\phi^n))$ is Cauchy in $L^2(\Omega)$.
\end{proposition}

\begin{proof}
By linearity,
$I_0(\phi^n)-I_0(\phi^m)=I_0(\phi^n-\phi^m)$, and apply It\^o isometry to the simple process $\phi^n-\phi^m$.
\end{proof}

\textbf{To do on board}
Key point: the integral is an \textbf{isometry}. Therefore the $L^2$ limit is natural.
\textbf{To do on board}

% ==========================================================
\section{Definition for general integrands (the $L^2$ limit)}
% ==========================================================

\begin{definition}[It\^o integral]
Let $\phi$ be adapted and satisfy $\E\!\int_0^T \phi_t^2\,\dd t<\infty$.
Choose simple processes $\phi^n$ such that $\phi^n\to\phi$ in $L^2(\Omega\times[0,T])$.
Define
\[
\int_0^T \phi_t\,\dd W_t := \lim_{n\to\infty} I_0(\phi^n)
\quad\text{in } L^2(\Omega).
\]
\end{definition}

\subsection{Well-definedness (independence of the approximation)}

\begin{theorem}[Well-definedness]
The limit above does not depend on the choice of the approximating sequence $(\phi^n)$.
\end{theorem}

\begin{proof}
Let $\phi^n\to\phi$ and $\psi^n\to\phi$ in $L^2(\Omega\times[0,T])$. Then $\phi^n-\psi^n\to 0$ in that space, and
by It\^o isometry for simple processes,
\[
\E\!\left[\big(I_0(\phi^n)-I_0(\psi^n)\big)^2\right]
=\E\!\left[\int_0^T (\phi_t^n-\psi_t^n)^2\,\dd t\right]\xrightarrow[n\to\infty]{}0.
\]
Hence $I_0(\phi^n)-I_0(\psi^n)\to 0$ in $L^2(\Omega)$, so the limits coincide.
\end{proof}

\textbf{To do on board}
Clean uniqueness proof: the $L^2$ metric + isometry does all the work.
\textbf{To do on board}

% ==========================================================
\section{Main properties: martingale, covariance, time continuity}
% ==========================================================

\subsection{Martingale}

For $t\le T$, define
\[
M_t := \int_0^t \phi_s\,\dd W_s.
\]

\begin{theorem}[Martingale property]
If $\phi$ is adapted and $\E\!\int_0^T \phi_t^2\,\dd t<\infty$, then $(M_t)_{0\le t\le T}$
is an $(\F_t)$-martingale and $\E[M_t]=0$.
\end{theorem}

\begin{proof}[Proof (simple case first, then extension)]
\textbf{Simple case.} Let $\phi_t=\sum_k \phi_{t_k}\1_{(t_k,t_{k+1}]}(t)$.
Fix $s<t$ and use a partition containing $s$. Then
\[
M_t-M_s=\sum_{k:\,t_k\in(s,t]} \phi_{t_k}\big(W_{t_{k+1}}-W_{t_k}\big).
\]
Condition on $\F_s$: each term has zero conditional mean (by iterated conditioning), hence
$\E[M_t-M_s\mid\F_s]=0$ and therefore $\E[M_t\mid\F_s]=M_s$.

\textbf{General case.} Approximate $\phi$ by simple $\phi^n$ in $L^2(\Omega\times[0,T])$.
Then $M_t^n:=\int_0^t \phi_s^n\,\dd W_s$ are martingales and, by It\^o isometry,
\[
\E[(M_t^n-M_t)^2]=\E\!\int_0^t(\phi_s^n-\phi_s)^2\,\dd s\to 0.
\]
Use $L^2$-continuity of conditional expectation to pass to the limit in
$\E[M_t^n\mid\F_s]=M_s^n$.
\end{proof}

\textbf{To do on board}
\begin{itemize}
  \item On the board: do the \textbf{simple} case only (clear and sufficient for intuition).
  \item For the extension: ``$L^2$ approximation + continuity of $\E[\cdot\mid\F_s]$''.
\end{itemize}
\textbf{To do on board}

\subsection{Isometry at time $t$ and polarization}

\begin{proposition}[Isometry at time $t$]
\[
\E[M_t^2]=\E\!\left[\int_0^t \phi_s^2\,\dd s\right].
\]
\end{proposition}

\begin{proposition}[Covariance / polarization]
If $\phi,\psi$ are square-integrable, then
\[
\E\!\left[\Big(\int_0^t \phi_s\,\dd W_s\Big)\Big(\int_0^t \psi_s\,\dd W_s\Big)\right]
=\E\!\left[\int_0^t \phi_s\psi_s\,\dd s\right].
\]
\end{proposition}

\begin{proof}
Use polarization:
\[
4\langle A,B\rangle = \|A+B\|^2-\|A-B\|^2,
\]
with $A=\int_0^t\phi\,\dd W$ and $B=\int_0^t\psi\,\dd W$, and apply the isometry to $\phi\pm\psi$.
\end{proof}

\subsection{Time continuity (brief)}

\begin{proposition}[$L^2$ continuity]
For $s<t$,
\[
\E[(M_t-M_s)^2]=\E\!\left[\int_s^t \phi_u^2\,\dd u\right]\xrightarrow[t\to s]{}0.
\]
\end{proposition}

\begin{remark}
One can then cite: an $L^2$-continuous martingale admits a continuous modification
(standard results via Doob/BDG inequalities). We do not reprove this here.
\end{remark}

% ==========================================================
\section{Why right-point sampling creates an extra ``drift''}
% ==========================================================

\subsection{Canonical example: $\phi_t=W_t$}

On a uniform partition:
\[
L_n=\sum_k W_{t_k}\Delta W_k,\qquad
R_n=\sum_k W_{t_{k+1}}\Delta W_k.
\]
We have $R_n=L_n+\sum_k(\Delta W_k)^2$, hence $R_n=L_n+T+o_{\Pbb}(1)$.

\begin{proposition}[Limit identification (formula to remember)]
\[
\int_0^T W_t\,\dd W_t=\frac12\big(W_T^2-T\big).
\]
\end{proposition}

\begin{proof}[Proof via sums]
Use the algebraic identity
\[
W_{t_{k+1}}^2-W_{t_k}^2
=2W_{t_k}\Delta W_k+(\Delta W_k)^2.
\]
Summing over $k$:
\[
W_T^2-W_0^2=2\sum_k W_{t_k}\Delta W_k+\sum_k(\Delta W_k)^2.
\]
Therefore
\[
L_n=\sum_k W_{t_k}\Delta W_k
=\frac12\Big(W_T^2-\sum_k(\Delta W_k)^2\Big)
\xrightarrow[n\to\infty]{\Pbb}\frac12(W_T^2-T),
\]
because $\sum_k(\Delta W_k)^2\to T$ in probability.
By definition, the limit of $L_n$ is $\int_0^T W_t\,\dd W_t$.
\end{proof}

\textbf{To do on board}
\begin{itemize}
  \item Very accessible proof: only algebra + quadratic variation.
  \item Point out: the term $-T/2$ is exactly the It\^o correction.
\end{itemize}
\textbf{To do on board}

% ==========================================================
\section{Deterministic integrand: Gaussianity and variance}
% ==========================================================

\begin{theorem}[Deterministic case]
If $f\in L^2(0,T)$ is deterministic, then
\[
\int_0^T f(t)\,\dd W_t \sim \mathcal N\!\Big(0,\ \int_0^T f(t)^2\,\dd t\Big).
\]
\end{theorem}

\begin{proof}[Clean and short proof idea]
Approximate $f$ by step functions $f^n=\sum_k c_k\1_{(t_k,t_{k+1}]}$ in $L^2(0,T)$.
Then
\[
\int_0^T f^n(t)\,\dd W_t=\sum_k c_k\big(W_{t_{k+1}}-W_{t_k}\big),
\]
a sum of independent Gaussian variables, hence Gaussian with mean $0$ and variance
$\sum_k c_k^2(t_{k+1}-t_k)=\int_0^T (f^n)^2$.
By It\^o isometry, $\int f^n\,\dd W\to \int f\,\dd W$ in $L^2$, hence in law, and the variances converge to $\int f^2$.
The limit of centered Gaussians whose variances converge is centered Gaussian with the limit variance.
\end{proof}

% ==========================================================
\section{Quadratic variation link: $[M]_t=\int_0^t \phi^2$ (preview)}
% ==========================================================

\begin{theorem}[Quadratic variation of an It\^o integral (statement)]
If $M_t=\int_0^t \phi_s\,\dd W_s$, then
\[
[M]_t=\int_0^t \phi_s^2\,\dd s.
\]
\end{theorem}

\begin{proof}[Sketch (to be detailed later)]
For a partition $\pi$ of $[0,t]$,
\[
\sum_k (M_{t_{k+1}}-M_{t_k})^2
=\sum_k \left(\int_{t_k}^{t_{k+1}}\phi_s\,\dd W_s\right)^2.
\]
Decompose
\[
\int_{t_k}^{t_{k+1}}\phi_s\,\dd W_s
=\phi_{t_k}(W_{t_{k+1}}-W_{t_k})+\int_{t_k}^{t_{k+1}}(\phi_s-\phi_{t_k})\,\dd W_s,
\]
and control the error by It\^o isometry:
\[
\E\left[\left(\int_{t_k}^{t_{k+1}}(\phi_s-\phi_{t_k})\,\dd W_s\right)^2\right]
=\E\int_{t_k}^{t_{k+1}}(\phi_s-\phi_{t_k})^2\,\dd s \to 0
\quad\text{as }|\pi|\to 0.
\]
The main term yields
\[
\sum_k \phi_{t_k}^2 (W_{t_{k+1}}-W_{t_k})^2 \approx \sum_k \phi_{t_k}^2 (t_{k+1}-t_k)
\to \int_0^t \phi_s^2\,\dd s.
\]
\end{proof}

\textbf{To do on board}
You can say: ``we will justify this properly when we develop covariation''.
But giving the formula now is useful: it explains the It\^o correction later.
\textbf{To do on board}

% ==========================================================
\section{Preview of the next lecture: toward It\^o's formula}
% ==========================================================

\begin{remark}[Roadmap]
Next lecture: covariation $[X,Y]$, the multiplication table ($dW^2=dt$),
product rule, then It\^o for $f(t,X_t)$.
Main takeaway today:
\[
\int_0^t \phi\,\dd W \text{ is a martingale and }
\E\left[\left(\int_0^t \phi\,\dd W\right)^2\right]=\E\int_0^t \phi^2.
\]
\end{remark}

\end{document}